\section{The five finite-index steps}\label{fc:section}

We now compose the comparisons at the actual integers
$n=26,\ldots,30$. The predecessor in each step is the unrestricted
bound $C_{n-1}\le\CE$. All parameters in a single step describe
one extremizer supplied by Lemma~\ref{gi:extremizer}.
In particular, neither its third moment nor its noise is chosen
independently for different count layers.
Throughout this section $R(t)$ denotes \eqref{pp:actual-residual}
with the fixed target $c_n=\CE/\sqrt n$. Computations can use
a smaller directed rational coefficient: an upper bound for that
larger residual is also an upper bound for this $R(t)$.

\subsection{Localization and the domain of the calculation}
The supporting localization proof gives three closed moment ranges.
For $\rho\ge6/5$, a support-free Prawitz comparison at each actual
integer gives $R_n(P)<.4097<\CE$; the large-moment ray is included.
For $1\le\rho\le11/10$, the contact and resource reductions,
with the full predecessor bound, give complete conditional count
covers. The original covers across the five integers consist of $79$ probability
intervals, $2289$ count roots and $74170$ terminal inequalities.
Every other probability interval is excluded already during
localization. These counts describe the coverage, not a replacement
for the inequalities proving it. The present composition uses those
ordinary count proofs on $72$ intervals.
Proposition~\ref{fc:lowprofiles} covers the remaining seven.

It remains to handle $11/10\le\rho\le6/5$.
At each actual integer, localization covers the $100$ original
closed intervals $[i/100,(i+1)/100]$, $0\le i<100$.
It excludes $89$ intervals and leaves exactly
\[
 I_i=[i/100,(i+1)/100],\qquad 32\le i\le42.
\]
The four outer intervals $I_{32},I_{33},I_{41},I_{42}$ have
complete ordinary count covers. At $n=26$, the same is true of
$I_{40}$. The new comparisons therefore concern $34$ parameter
profiles: $I_{34},\ldots,I_{39}$ at $26$, and
$I_{34},\ldots,I_{40}$ at each of $27,28,29,30$.
Their scalar caps are listed in Appendix~\ref{app:parameters}.

A profile supplies inward and outward support caps
$\Delta,O_0,O_1$, a noise cap $\Lambda$, and the necessary joint
inequalities
\begin{equation}\label{fc:resources}
 a_0+a_1\le\Delta(1-\lambda/\Lambda),\qquad
g(a_0+a_1)+\lambda
       \{(1-\alpha)\mathfrak h_-+\alpha\mathfrak h_+\}
       \le\mathfrak R.
\end{equation}
The second inequality is the consequence of
\eqref{loc:affine-resource} with $D\le\overline D$ and
$\mathfrak R=\overline D-b$. It can strengthen the first cap and supplies
the signed plane \eqref{mc:signed-plane}. Its coefficients come
from the actual-index resource proof. They are held fixed on
their stated probability interval, not fitted to the count being
estimated. The contact bounds also give
\begin{equation}\label{fc:contact}
 -H_-\sqrt{pq+\lambda/n}
 \le k+t-np\le H_+\sqrt{pq+\lambda/n},\qquad |t|\le1/2.
\end{equation}
This leaves only three neighboring counts in each profile.
The support caps, both resource bounds and the actual-moment floor
are used simultaneously, with no finite-support assumption on the
original conditional laws.

\subsection{Bounding a whole parameter region}
For a current-count price, let $L_*$ be the acquired actual-noise cap.
The seven original closed noise intervals have endpoints
\[
 0,\quad\Lambda/8,\quad\Lambda/4,\quad5\Lambda/16,\quad
 3\Lambda/8,\quad7\Lambda/16,\quad\Lambda/2,\quad L_*;
\]
the four allocation intervals have endpoints $0,1/4,1/2,3/4,1$.
These endpoints are increasing in every profile.
Their Cartesian product gives $28$ fixed reference laws per
current-count family. Each is valid for its entire closed region
by Lemma~\ref{mc:two-moment}; it is not a sample at its midpoint.
The signed plane and the inward cap restrict the admissible
conditional variances before the reference law is formed.
The noise correction in Proposition~\ref{ce:fixed-box} preserves
the difference between that reference law and the actual noise.
Probability faces and phase faces extend to the whole region
by the corresponding lower-curvature bounds. A tilted phase
face at the Gaussian inflection is used when needed.

For a reflected-square price, Proposition~\ref{ce:coherent}
keeps the actual noise and allocation until the two derivative
gates have been proved. The resulting endpoint envelope uses
the same moment floor as the original law. A positive-atom
activation is covered by the weak derivatives
\eqref{pp:activation-curvature}. Complementary prices may then be
combined through the finite lower envelope of Lemma~\ref{pp:mixture},
using their directed lower target coefficients. Only explicit reference prices are
differentiated in these steps.

All parameter endpoints, chord coefficients and interval enclosures
are rational. The square roots in the finite prices, and the
normal CDF and density, are enclosed outwards. A stationary
width is accepted only when it belongs to its asserted active
set. Curvature bounds include activation boundaries. The finite
maximum of every interpolating quadratic includes its interior
vertex when present. Thus a negative upper bound applies to
every point of the closed domain, including all shared faces.
The full count roots and the unions of their closed phase
intervals are retained when these regional bounds are composed.

\subsection{Examples of the finite comparisons}
An illustrative replacement occurs at $n=26$, $p\in I_{35}$,
count $9$. The current-count price with the least legal slope
is strict on $[-1/8,9/32]$ except for the last phase strip of
the first noise interval. On that first noise interval,
$[0,\Lambda/8]$, retain the current price through $t=1/4$ and
use a coherent count-$10$ reflected-square price of width $9/8$
on $[1/4,9/32]$. Its allocation and noise derivative margins
are respectively at least $.120534962164$ and $.072056354791$,
and its variance derivative is at least $.759364977755$.
Its residual upper bound is $-.018464965470$ for its directed
lower target coefficient. The other original noise intervals
retain the current price on the whole strip.
The resulting uniform upper bound on $[-1/8,9/32]$ is
\begin{equation}\label{fc:n26-simple}
 R(t)\le-.000067780727.
\end{equation}
The adjacent negative phases have a full-square or current-count
bound, and $[9/32,1/2]$ has the full-square count-$10$ bound.
Consequently this entire count layer uses only second-moment
comparisons. The remaining count-$10$ regions in $I_{36},I_{37}$
likewise admit current-count comparisons with the rational slope
multipliers $1,5/4,3/2$. No fourth-moment price or optimized
deterministic-group slope is needed by these analytic profile comparisons.

A second completion illustrates why it is useful to distinguish
the actual CDF from its price. At $n=30$, $p\in I_{37}$,
the count-$10$ contact interval begins at the exact outward
bound
\[
 a_{37}=249467754491/10^{12}<1/4.
\]
An endpoint estimate at $1/4$ extends over $[a_{37},1/4]$
by Lemma~\ref{pp:cdf-bridge}. The comparison used to prove that
endpoint estimate need not remain valid on the short interval.
Monotonicity of the original CDF is sufficient.

\begin{proposition}[Seven lower-moment profiles]\label{fc:lowprofiles}
For each row below, retain the original actual-$n$ support and resource
acquisition and the unrestricted predecessor $C_{n-1}\le\CE$.
The whole profile $p\in I_i$, $1\le\rho\le11/10$ has conditional
count residual bounded by the last column. The joint intercept
$\Lambda$ is retained; $L_*$ is a necessary ceiling on the actual noise.
\begin{center}\small\setlength{\tabcolsep}{3pt}
\begin{tabular}{ccrrrr}\toprule
$n$ & $I_i$ & $\Delta$ & $\Lambda$ & $L_*$ & residual upper\\\midrule
26 & $I_{39}$ & .05769 & .30708 & .17375152 & $-.000098710877$\\
26 & $I_{40}$ & .057809 & .301161 & .17805629 & $-.000005610547$\\
26 & $I_{41}$ & .054236 & .276951 & .162619 & $-.000049410586$\\
27 & $I_{41}$ & .052227 & .276951 & .16129388 & $-.000016889$\\
28 & $I_{41}$ & .050362 & .276951 & .16009718 & $-.000027176644$\\
29 & $I_{41}$ & .048625 & .276951 & .15901083 & $-.000059169468$\\
30 & $I_{41}$ & .047004 & .276951 & .15802002 & $-.000017416596$\\
\bottomrule\end{tabular}\end{center}
\end{proposition}
\begin{proof}
Appendix~\ref{app:finite-details} gives all seven finite plans and their
closed count/phase unions. Each uses the seven original noise bands and
four allocation quarters. The signed resource plane first gives $L_*$;
current-count and reflected-square prices then cover every surviving
contact interval, with the bounds in the table.

Three features require the full comparison rules: reversed signed
prices use \eqref{loc:two-endpoint-noise}; the count-$12$ reference at
$n=29$ crosses $np$ and uses both branches of
\eqref{ce:variance-endpoint-max}; and the early $n=30$ square has
$y_{\max}>0$, covered by the $y_{\max}<b$ bound in Section~\ref{pp:section}.
The remaining steps use the same endpoint, curvature and threshold-shift
identities, with each row's actual integer and acquired parameters.
\end{proof}

\subsection{The finite implications}


The complete middle-moment profiles at $27$--$30$ give the following outward
upper bounds on $R(t)$, uniformly over their probability interval,
actual noise, allocation, all surviving counts and contact phases.
For readability, the table reports $10^6$ times these bounds,
rounded upwards to three decimal places.
\begin{center}
\begin{tabular}{crrrr}
\toprule
interval & $n=27$ & $n=28$ & $n=29$ & $n=30$\\
\midrule
$I_{34}$ & $-320.741$ & $-263.994$ & $-321.528$ & $-241.967$\\
$I_{35}$ & $-297.200$ & $-134.578$ & $-287.407$ & $-436.298$\\
$I_{36}$ & $-294.861$ & $-12.020$  & $-96.700$  & $-156.955$\\
$I_{37}$ & $-145.182$ & $-12.884$  & $-1.810$   & $-63.857$\\
$I_{38}$ & $-70.159$  & $-56.159$  & $-79.848$  & $-22.937$\\
$I_{39}$ & $-35.366$  & $-4.435$   & $-585.317$ & $-75.052$\\
$I_{40}$ & $-1160.650$& $-723.488$ & $-713.730$ & $-987.882$\\
\bottomrule
\end{tabular}
\end{center}
The six analytic profiles at $26$ are also strictly negative;
\eqref{fc:n26-simple} describes its only former fourth-moment
region. The full results include the exact resource caps, count
domains, price choices, derivative gates and outward rational
bounds, rather than only these decimal margins.

\begin{proposition}[The five actual-index implications]\label{fc:five}
For each $n=26,27,28,29,30$,
\[
 C_{n-1}\le\CE\quad\Longrightarrow\quad C_n\le\CE.
\]
\end{proposition}
\begin{proof}
If $C_n>\CE$, Lemma~\ref{gi:extremizer} supplies a positive
closed extremizer. The full predecessor bound makes its
allowance \eqref{gi:predecessor} zero, so the actual-index
localization and resource inequalities apply to this same law.
The support-free Fourier estimate excludes $\rho\ge6/5$.
The complete lower-moment count covers exclude
$1\le\rho\le11/10$.
In the remaining closed band, localization either gives an
immediate contradiction or places the law in one of the eleven
intervals above. The ordinary count covers and the $34$ analytic
profiles exclude them all. Every surviving count and its entire
contact interval has a strict negative bound, contradicting the
positive extremal score. The three closed moment ranges exhaust
$\rho\ge1$.
\end{proof}

\begin{theorem}[The iid Esseen-constant inequality]\label{fc:main}
For every standardized real distribution $P$ with
$\rho(P)=\E|X|^3<\infty$ and every positive integer $n$,
\[
 \sqrt n\,\Delta_n(P)\le
 \frac{3+\sqrt{10}}{6\sqrt{2\pi}}\,\rho(P).
\]
The constant is optimal over this class and all $n$.
\end{theorem}
\begin{proof}
Input~(I) of Proposition~\ref{gi:inputs} supplies the unrestricted
bound at each index through $25$. Apply Proposition~\ref{fc:five}
in order, discharging each global predecessor, to obtain the
bounds through $30$. In particular $C_{30}<\CE+1/100$, so
input~(II) gives $31$--$46$. Proposition~\ref{gi:tail} then gives
every $n\ge47$ from its two sharp predecessors and the uniform
local exclusion (III).
These four ranges exhaust all positive integers.

For the standardized Bernoulli law with
$p_*=(4-\sqrt{10})/2$, the classical binomial asymptotic gives
\[
 \lim_{n\to\infty}\frac{\sqrt n\,\Delta_n(P_*)}{\rho(P_*)}=\CE;
\]
see Schulz \cite[Remarks~4.1 and 4.2]{Schulz2016}.
Thus a smaller uniform constant is impossible.
\end{proof}

\section{Scope of the supporting material}\label{fc:supporting}
The finite comparisons in Sections~\ref{mc:section}--\ref{ce:section}
are analytic inequalities for arbitrary real conditional laws.
The finite calculations certify their use on the closed parameter
domains of Section~\ref{fc:section}. The supporting material also
contains the continuous localization and the proofs of inputs
(I)--(III). These are essential parts of the all-index argument,
not hypotheses inferred from the numerical table.

The moment-defect and localization framework, the original
subdivisions, and the imported count and Fourier arguments are
due to Rodziewicz and Ko\l{}odziejek \cite{RK2026}.
The role of the present comparison method is to close the remaining
general-law regions at the five actual integers and compose the
supporting implications in a well-founded order.

This working version includes the companion mathematical sources,
exact acquired parameters, collected frozen arithmetic inputs and
archived calculation outputs. The accompanying programs locate these
materials and run their mathematical readers from the collected
directories. Fresh numerical regeneration uses the retained producer
entrypoints and their stated runtime requirements. The collection is
local at this stage; no permanent public archive identifier is assigned
in this draft.
